% -----------------------------
% 執筆にあたっての留意点
% -----------------------------
%
% おわりに（結論）の章は，論文を読み終えた読者の頭に何を残すかを決める章です．
% 分量は多くありませんが，アブストラクトと並んで「そこだけ読む読者」が最も多い章でもあります．
% 参考：how-to-xx/how-to-write-a-paper/1-学術論文の書き方.md 3.6節
%
% 以下の順に書くとよいでしょう．
%
% 1. 本研究で何をしたか（1〜2文）
% 2. 何がわかったか（定量的な結果を使って具体的に）
% 3. その知見が持つ意味
% 4. 今後の課題（Future Work）
%
% ■ アブストラクトやイントロダクションのコピーにしてはいけません
% 最もやりがちな失敗です．結論の章は，イントロダクションで述べた主張を「実験を終えた今の言葉」で
% 述べ直す場です．たとえばイントロダクションで「検索行動を大きく変える」と書いたのであれば，
% 結論では「セッションあたりの検索回数が◯倍に増加することを確認した」と，
% 得られた数値を使って具体的に言い直します．
%
% ■ ここで新しい主張を初めて出してはいけません
% 結論に書いてよいのは，本文で根拠を示した内容だけです．
% 「本研究の結果は〜にも応用できるだろう」と書きたくなったら，それが結果に支えられているかを確認し，
% 支えられていないなら今後の課題として書いてください．
%
% ■ 今後の課題は，研究の価値の一部です
% 「〜が今後の課題である」と列挙して終わるのではなく，本研究の限界のうちどれが次に解くべき問題なのかを，
% 理由とともに述べてください．すでに取り組み始めている研究があれば，
% 「現在〜に取り組んでおり，予備的な結果は有望である」と書いておくと，
% そのテーマが自分のものであることを示せます．
%
% ■ 「今後の課題」が言い訳の場にならないようにしてください
% 「時間がなかったため〜は実施できなかった」といった書き方は，読者にとって何の情報にもなりません．
% 「〜を明らかにするには〜の実験が必要であり，これを次の課題とする」と，次の研究の設計として書きます．

% ----------------------
% おわりに/conclusion
% ----------------------
\section{おわりに}
